% part: intuitionistic-logic
% chapter: propositions-as-types
% section: proof-terms

\documentclass[../../../include/open-logic-section]{subfiles}

\begin{document}

\olfileid{int}{pty}{ter}

\olsection{સાબિતી-પદો}

\Log{N1i}માં ધારણાઓને મુક્તિનાં ચિહ્ન તરીકે
ચલો આપીએ છીએ (દાખલા તરીકે $!A^x$).
\Log{N2i}માં દરેક સિક્વન્ટ $\Gamma \Sequent !A$ની
ડાબી બાજુના સંદર્ભ~$\Gamma$માં ચલ--!!{formula}
જોડીઓની યાદી છે. તેથી \Log{N2i}માં
!!a{proof}નો દરેક સિક્વન્ટ એટલું જ નથી
જણાવતો કે અત્યાર સુધીની !!{proof}એ શું
!!{prove}d કર્યું છે (એટલે $!A$), પણ દરેક
સ્થાને કઈ ધારણાઓ ખુલ્લી છે (એટલે $\Gamma$)
તે પણ જણાવે છે. દરેક સિક્વન્ટમાં~$!A$
\emph{કેવી રીતે} !!{prove}d થયું છે તે
માહિતી પણ ઉમેરીએ તો? આ માટે
\emph{પદથી ચિહ્નિત} તંત્ર \Log{tN2i}
વ્યાખ્યાયિત કરી શકીએ છીએ, જેમાં સિક્વન્ટનું
સ્વરૂપ $\Gamma \Sequent N : !A$ છે.
$\Gamma$ અને $!A$ પહેલાં જેવા જ છે:
$x: !B$ સ્વરૂપની જોડીઓનો ગણ અને
સિક્વન્ટ પર પૂરી થતી ઉપ-!!{proof}એ
સંદર્ભ~$\Gamma$ની ધારણાઓમાંથી અનુસરતું
બતાવેલું !!{formula}. $N$ એ આ સિક્વન્ટ
પર પૂરી થતી !!{proof}ને સંકેતરૂપે
રજૂ કરતું ચોક્કસ પદ હશે.

દરેક સિક્વન્ટને સોંપેલા પદો $x$, $y$,
\dots, ચલોથી બને છે. તેમાં અત્યાર સુધીનાં
અનુમાનો સંકેતરૂપે રજૂ કરતાં વિધેય-ચિહ્નો
વાપરીએ છીએ. દરેક અનુમાન-નિયમ માટે અલગ
વિધેય-ચિહ્ન જોઈએ. !!{introduction}
નિયમોને અનુરૂપ ચિહ્નો ``રચકો'' અને
!!{elimination} નિયમોને અનુરૂપ ચિહ્નો
``વિઘટકો'' કહેવાય છે. નિયમનાં જેટલાં
આધારવિધાનો હોય, દરેક ચિહ્ન એટલાં
આર્ગ્યુમેન્ટ લે છે. દાખલા તરીકે,
\Intro{\land} અને \Elim{\land}[1]નાં
ચિહ્નો $c^\land$ (બે-સ્થાની) અને
$d^\land_1$ (એક-સ્થાની) હોઈ શકે.
\Log{tN2i}માં આ નિયમો આમ થાય છે:
\[
\Axiom$\Gamma_1 \fCenter N: !A$
\Axiom$\Gamma_2 \fCenter M: !B$
\RightLabel{\Intro{\land}}
\BinaryInf$\Gamma_1, \Gamma_2 \fCenter c^\land(N, M) : !A \land !B $
\DisplayProof
\quad
\Axiom$\Gamma \fCenter N: !A \land !B$
\RightLabel{\Elim{\land}[1]}
\UnaryInf$\Gamma \fCenter d^\land_1(N) : !A$
\DisplayProof
\]
\sourcecorrection{OLMD-287}{મૂળ સંયોજનના રચકમાં નાનું $m$ છે; બીજા આધારવિધાનનું પદ $M$ હોવાથી અહીં મોટું $M$ કર્યું છે.}
અનુમાન-નિયમ ધારણા મુક્ત કરે, એટલે તેના
આધારવિધાનમાં ચિહ્નિત !!{formula}~$x:!A$
હોય પણ નિષ્કર્ષના સંદર્ભમાં ન હોય, તો
સાબિતી-પદમાં પણ આ બતાવવું જોઈએ.
આવા નિયમોનાં વિધેય-ચિહ્નો સંબંધિત
ચલોને \emph{બદ્ધ કરે} છે. દાખલા તરીકે,
\Intro{\lif} નિયમ આધારવિધાન
$x:!A, \Gamma \Sequent !B$માંથી
$\Gamma \Sequent !A \lif !B$ આપે છે.
રચક $c^{\lif}$ એક આર્ગ્યુમેન્ટ લે છે,
પણ ચલ~$x$ બદ્ધ થાય છે તે પણ
બતાવવું પડે. સાબિતી-પદ આ રીતે
ચિહ્ન~$x$વાળી ધારણા મુક્ત થઈ છે
તે વ્યક્ત કરે છે. દાખલા તરીકે,
આધારવિધાનનું સાબિતી-પદ~$N$ હોય,
તો નિષ્કર્ષનું પદ $c^{\lif}([x]N)$
અને નિયમ આમ લખી શકીએ:
\[
\Axiom$x:!A, \Gamma \fCenter N: !B$
\RightLabel{\Intro{\lif}}
\UnaryInf$\Gamma \fCenter c^{\lif}([x]N) : !A \lif !B$
\DisplayProof
\]

!!a{proof}ને સંપૂર્ણ રીતે ફરી રચવા માટે
વધારાની માહિતી જોઈએ. નિયમના આધારવિધાનનાં
!!{formula}sથી નિષ્કર્ષનું !!{formula}
સંપૂર્ણ નક્કી ન થતું હોય, તો તે માહિતી
પદમાં ઉમેરવી પડે. \Intro{\lif}માં આવું
થાય છે, એટલે $c^{\lif}([x:!A]N)$
લખીશું. \Intro{\lor} અને~\FalseIntમાં
પણ આવું થાય છે, એટલે એ માહિતી ધરાવતાં
વિધેય-ચિહ્નો વાપરીએ છીએ. દાખલા તરીકે,
\[
\Axiom$\Gamma \fCenter N:!A$
\RightLabel{\Intro{\lor}[1]}
\UnaryInf$\Gamma \fCenter c^{\lor{!B}}_1(N):!A \lor !B$
\DisplayProof
\quad
\Axiom$\Gamma \fCenter N:\lfalse$
\RightLabel{\FalseInt}
\UnaryInf$\Gamma \fCenter d^\lfalse_{!A}(N):!A$
\DisplayProof
\]
\sourcecorrection{OLMD-288}{મૂળ વિકલ્પ-પરિચયના નિષ્કર્ષમાં રચકનો આર્ગ્યુમેન્ટ છૂટી ગયો છે; આધારવિધાનનું પદ $N$ હોવાથી તેને ઉમેર્યું છે.}

આ વિચારો પરથી સિક્વન્ટ-શૈલીનું સ્વાભાવિક
નિગમન તંત્ર રચી શકાય છે, જેમાં દરેક
સિક્વન્ટનું સ્વરૂપ $\Gamma \Sequent N : !A$
છે અને~$N$ એક \emph{સાબિતી-પદ} છે.

આવાં સાબિતી-પદો અને \emph{પ્રકારિત લૅમ્બડા-કલન}નાં
પદો વચ્ચે નજીકનો સંબંધ છે. આ સંબંધને
અનુસરવા ઉપરનાં સામાન્ય રચક અને વિઘટકનાં
ચિહ્નોની જગ્યાએ લૅમ્બડા-કલનના સાહિત્યમાં
વપરાતાં ચિહ્નો લઈશું. દાખલા તરીકે,
$c^{\lif}([x:!A]N)$ની જગ્યાએ
$\lambd[\typeof{x}{!A}][N]$,
$d^{\lif}(N,M)$ની જગ્યાએ $(NM)$,
$c^\land(N,M)$ની જગ્યાએ $\pair{N}{M}$,
અને~$d^\land_i(N)$ની જગ્યાએ $\proj{i}{N}$
લખીએ.

\begin{defn}
  સાબિતી-પદો આ નિયમો વડે આગમનાત્મક રીતે બને છે:
  \begin{enumerate}
  \item દરેક ચલ~$x$ સાબિતી-પદ છે (સ્વયંસિદ્ધો).
  \item $x$ ચલ હોય, $!A$ !!a{formula} હોય અને
    $N$ સાબિતી-પદ હોય, તો $\lambd[\typeof{x}{!A}][N]$
    પણ સાબિતી-પદ છે~(\Intro\lif).
  \item $N$ અને $M$ સાબિતી-પદો હોય, તો
    $(NM)$ પણ સાબિતી-પદ છે~(\Elim\lif).
  \item $N$ અને $M$ સાબિતી-પદો હોય, તો
    $\pair{N}{M}$ સાબિતી-પદ છે~(\Intro\land).
  \item $N$ સાબિતી-પદ હોય, તો $\proj{i}{N}$
    પણ સાબિતી-પદ છે; અહીં $i$ એ $1$ અથવા~$2$
    છે~(\Elim\land[i]).
  \item $N$ સાબિતી-પદ હોય અને $!A$ સૂત્ર હોય,
    તો $\inj[!A]{i}{N}$ સાબિતી-પદ છે;
    અહીં $i$ એ $1$ અથવા~$2$ છે~(\Intro\lor[i]).
    \sourcecorrection{OLMD-289}{મૂળ આ કિસ્સામાં સાબિતી-પદ $N$થી શરૂ કરે છે, પણ નિષ્કર્ષમાં અસંબંધિત $!M$ આવે છે. અહીં નિષ્કર્ષનું પદ $N$ કર્યું છે.}
  \item $M$, $N_1$, $N_2$ સાબિતી-પદો હોય
    અને $x_1$, $x_2$ ચલો હોય, તો
    $\dcase{M}{x_1}{N_1}{x_2}{N_2}$
    સાબિતી-પદ છે~(\Elim\lor).
  \item $N$ સાબિતી-પદ હોય અને $!A$
    !!a{formula} હોય, તો $\abort{!A}{N}$
    સાબિતી-પદ છે~(\FalseInt).
  \end{enumerate}
\end{defn}

દરેક સાબિતી-પદ કોઈ !!a{proof}નું
સાબિતી-પદ હોતું નથી. દાખલા તરીકે,
$\pair{N}{M}$ પદ \Intro{\land}
અનુમાનથી પૂરી થતી સાબિતીને સંકેતરૂપે
રજૂ કરે છે, એટલે તેનું અંતિમ !!{formula}
સંયોજન~$!A \land !B$ છે. તે
\Elim{\lif} નિયમનું મુખ્ય આધારવિધાન
ન હોઈ શકે, એટલે $(\pair{N}{M}P)$
યોગ્ય સાબિતીનું પદ ન હોઈ શકે.
\Log{tN2ip}ના નિયમો અનુસરતાં
સાબિતી-પદો જ \emph{યોગ્ય} છે.
આ નિયમો \olref{tab:tN2ip}માં આપ્યા છે.

\subfile{rules-tN2}

\end{document}
